%======================================================================
\section{Implicational intermediate logics}\label{sec:implicational}

In this section we construct, for every implicational formula $\Phi\in\CL$,
a formula $\Mz{\Phi}$ \kl{minimal in} $\CL$ with
$\IL\lplus\Mz{\Phi}=\IL\lplus\Phi$ (Theorem~\ref{thm:first-layer}). Its
\kl{order} is reduced to $4$ in Section~\ref{sec:order}.

By Proposition~\ref{prop:criterion}, we must make every \kl{splitting} and
every \kl{abstraction} of the output fail classically, while keeping a
substitution that recovers $\Phi$. A falsifying valuation for a modified
formula is easy to find when the modification is \emph{local}, that is,
when it affects a small \kl{read-once} piece of the formula on which the
rest of the formula imposes no condition. We achieve this by composing two
known translations with two general constructions:
\[
\Mz{\Phi}=\Tc{\Mon{\Cons_\Phi}{\nv{\rt}}} .
\]
Here $\Cons_\Phi$ is the list of \kl{local constraints} of $\Phi$, obtained
by \kl{renaming apart} the \kl{occurrences} of variables and applying the
polarized extension-variable translation (Section~\ref{sec:constraints}).
The \kl{monitor} $\Mon{\Cons}{r}$ turns any list of \kl{read-once} formulas
into a formula no \kl{separation} of which is classically valid
(Section~\ref{sec:localization}), and the \kl{tree closure} $\Tc{F}$ turns
any such formula into a \kl{minimal} one (Section~\ref{sec:minimality}).
The two general constructions do not depend on $\Phi$. We use Peirce's law
\[
\axP=((a\to b)\to a)\to a
\]
as a running example, and we explain after each step which
\kl{elementary generalizations} survive it.

\subsection{Two known translations}\label{sec:constraints}

\AP The first translation is \intro{renaming apart} the \kl{occurrences} of
variables. Let $\alpha$ be a formula. For every \kl{leaf} $o$ of $\alpha$
we take a new variable $\intro*\port{o}$, called a \intro{port}, distinct
\kl{leaves} receiving distinct \kl{ports}, and let $\intro*\ported{\alpha}$
be the result of replacing each \kl{leaf} $o$ by $\port{o}$; thus
$\ported{\alpha}$ is \kl{read-once}. If $a$ is the variable at $o$, the
\intro{guard} of $o$ is
\[
\intro*\guard{o}=
\begin{cases}
a\to\port{o}, & \text{if $o$ is \kl{positive} in $\alpha$},\\
\port{o}\to a, & \text{if $o$ is \kl{negative} in $\alpha$}.
\end{cases}
\]
\AP We write $\intro*\Guards_\alpha$ for the list of the \kl{guards} of
$\alpha$, call $a$ the \intro{center} of the \kl{ports} of its
\kl{occurrences}, and let $\intro*\collsub_\alpha$ be the substitution that
maps every \kl{port} to its \kl{center} and fixes all other variables. The
following is the usual monotonicity argument; the \kl{guards} are the
clauses for \kl{leaves} in the structure-preserving translation of Plaisted
and Greenbaum~\cite{PG86}.

\begin{lemma}\label{lem:port}
$\Guards_\alpha\vdash_{\IL}\alpha\to\ported{\alpha}$, and $\collsub_\alpha$
maps every \kl{guard} to a formula $a\to a$ and $\ported{\alpha}$ to
$\alpha$.
\end{lemma}
\begin{proof}
By induction on subformulas, $\Guards_\alpha$ derives
$\subf{\alpha}{u}\to\subf{\ported{\alpha}}{u}$ for \kl{positive} $u$ and
$\subf{\ported{\alpha}}{u}\to\subf{\alpha}{u}$ for \kl{negative} $u$, since
the antecedent of an implication has the opposite \kl{polarity}.
\end{proof}

\AP The second translation is the polarized form of Statman's
extension-variable translation~\cite{Sta79,PG86,Min90}. Let $\alpha$ be an
implicational formula. For each \kl{non-leaf occurrence} $u$ of $\alpha$ we
choose a new variable $\intro*\ev{u}$, called an \intro{extension variable},
distinct \kl{occurrences} receiving distinct variables. For every
\kl{occurrence} $u$ let $\intro*\rv{u}=\subf{\alpha}{u}$ if $u$ is a
\kl{leaf} and $\rv{u}=\ev{u}$ otherwise. \AP For a \kl{non-leaf
occurrence} $u$ put $\intro*\body{u}=\rv{u0}\to\rv{u1}$ and
\[
\intro*\dprem{u}=
\begin{cases}
\body{u}\to\ev{u}, & \text{if $u$ is \kl{positive}},\\
\ev{u}\to\body{u}, & \text{if $u$ is \kl{negative}}.
\end{cases}
\]

\begin{definition}\label{def:S}
\AP Let $u_1,\dots,u_m$ be the \kl{non-leaf occurrences} of $\alpha$ in
post-order; that is, $u_i$ precedes $u_j$ if and only if $u_i$ is
\kl{strictly below} $u_j$, or $u_i=w0u'$ and $u_j=w1u''$ for some words
$w,u',u''$. The \intro{translation} of $\alpha$ is
\[
\intro*\St(\alpha)=\dprem{u_1}\to\dots\to\dprem{u_m}\to\rv{\rt} .
\]
If $\alpha$ is a variable, then $\St(\alpha)=\alpha$.
\end{definition}

Every $\body{u}$ has \kl{order} $2$, so every $\dprem{u}$ has \kl{order} at
most $3$, and
\begin{equation}\label{eq:S-order}
\Ord(\St(\alpha))\le4,\qquad \sizeof{\St(\alpha)}=3\sizeof{\alpha}-2 .
\end{equation}
\AP Let $\intro*\decsub{\alpha}$ be the substitution that maps each
$\ev{u}$ to $\subf{\alpha}{u}$ and fixes all other variables. Then
$\decsub{\alpha}(\rv{u})=\subf{\alpha}{u}$, and $\decsub{\alpha}$ maps
every $\dprem{u}$ to $\subf{\alpha}{u}\to\subf{\alpha}{u}$.

\begin{lemma}\label{lem:S-sound}
For every implicational formula $\alpha$, the formulas $\dprem{u}$ derive
$\alpha\to\rv{\rt}$ in $\IL$; hence $\IL\vdash\alpha\to\St(\alpha)$.
Moreover $\IL\vdash\decsub{\alpha}(\St(\alpha))\to\alpha$, and therefore
$\IL\lplus\St(\alpha)=\IL\lplus\alpha$.
\end{lemma}
\begin{proof}
By induction on subformulas, the formulas $\dprem{u}$ derive
$\subf{\alpha}{u}\to\rv{u}$ for \kl{positive} $u$ and
$\rv{u}\to\subf{\alpha}{u}$ for \kl{negative} $u$. For instance, if $u$ is
\kl{positive}, then from $\rv{u0}\to\subf{\alpha}{u0}$ and
$\subf{\alpha}{u1}\to\rv{u1}$ we obtain
$\subf{\alpha}{u}\to\body{u}$, and $\dprem{u}=\body{u}\to\ev{u}$ gives
$\subf{\alpha}{u}\to\ev{u}$. The second claim holds since
$\decsub{\alpha}(\St(\alpha))=(\gamma_1\to\gamma_1)\to\dots\to
(\gamma_m\to\gamma_m)\to\alpha$.
\end{proof}

Every formula $\dprem{u}$ is used exactly once in this derivation, which
in fact takes place in $\BCI$.

\AP We combine the two translations. Let $\Phi$ be an implicational formula
that is not a variable. For every \kl{occurrence} $v$ of $\Phi$ we write
$\intro*\nv{v}$ for the \kl{port} of $v$ if $v$ is a \kl{leaf}, and for the
\kl{extension variable} of $v$ in $\ported{\Phi}$ otherwise; these are the
\intro{node variables}. \AP The list $\intro*\Cons_\Phi$ of \intro{local
constraints} of $\Phi$ consists of the \kl{guards} of $\Phi$ and the
formulas $\dprem{u}$ of $\ported{\Phi}$. Explicitly, writing
$\intro*\cons{v}$ for the constraint of $v$, with $u=v0$ and $w=v1$,
\[
\begin{array}{lll}
\subf{\Phi}{v} & \text{$v$ \kl{positive}} & \text{$v$ \kl{negative}}\\\hline
\text{a variable }p & p\to\nv{v} & \nv{v}\to p\\
\subf{\Phi}{u}\to\subf{\Phi}{w} & (\nv{u}\to\nv{w})\to\nv{v} & \nv{v}\to(\nv{u}\to\nv{w})
\end{array}
\]
There are $2\sizeof{\Phi}-1$ \kl{local constraints}; each is \kl{read-once},
has two or three variables, and has \kl{order} at most $3$. By
Lemmas~\ref{lem:port} and~\ref{lem:S-sound}:

\begin{lemma}\label{lem:constraints}
$\Cons_\Phi\vdash_{\IL}\Phi\to\nv{\rt}$. The substitution
$\collsub_\Phi\circ\decsub{\ported{\Phi}}$, which maps $\nv{v}$ to
$\subf{\Phi}{v}$ and fixes the variables of $\Phi$, maps every \kl{local
constraint} to a formula $\gamma\to\gamma$ and $\nv{\rt}$ to $\Phi$.
\end{lemma}

\begin{example}\label{ex:P-constraints}
The \kl{occurrences} of $\axP$, their \kl{polarities} and their \kl{local
constraints} are as follows.
\[
\begin{array}{llll}
v & \subf{\axP}{v} & & \cons{v}\\\hline
\rt & ((a\to b)\to a)\to a & + & (\nv{0}\to\nv{1})\to\nv{\rt}\\
0 & (a\to b)\to a & - & \nv{0}\to\nv{00}\to\nv{01}\\
00 & a\to b & + & (\nv{000}\to\nv{001})\to\nv{00}\\
000 & a & - & \nv{000}\to a\\
001 & b & + & b\to\nv{001}\\
01 & a & - & \nv{01}\to a\\
1 & a & + & a\to\nv{1}
\end{array}
\]
Each constraint says that a \kl{node variable} is an upper bound (for
\kl{positive} nodes) or a lower bound (for \kl{negative} nodes) of the
subformula at its node, in terms of the \kl{node variables} of the children.
\end{example}

The formula $\Cons_\Phi\to\nv{\rt}$, that is, $\St(\ported{\Phi})$ preceded
by the \kl{guards} of $\Phi$, generates the same \kl{logic} as $\Phi$ over
$\IL$ by Lemma~\ref{lem:constraints}, but it need not be \kl{minimal}, even
if $\Phi$ is. For $\axP$, the remaining constraints yield $\nv{\rt}$
classically without $b\to\nv{001}$: if $\nv{\rt}=0$, then $\nv{0}=1$ and
$\nv{1}=0$, hence $a=\nv{01}=\nv{000}=0$, hence $\nv{00}=1$, contradicting
$\nv{0}\to\nv{00}\to\nv{01}$. Therefore the \kl{separation} that replaces
$\nv{001}$ in $b\to\nv{001}$ by a new variable is still in $\CL$. The
problem is that a falsifying valuation must satisfy all the other
constraints at once, and these together express the validity of $\Phi$.

\subsection{The monitor}\label{sec:localization}

\AP Let $\Cons=(C_1,\dots,C_n)$ be a list of implicational formulas and $r$
a variable. We take new variables $\intro*\sel{1},\dots,\sel{n}$, called
\intro{selectors}, a new variable $\intro*\sw$, called the \intro{switch},
and a new variable $\intro*\goal$, and put
\[
\intro*\guarded{i}=\sel{i}\to C_i,\qquad
\intro*\link{i}=\sw\to\sel{i},\qquad
\intro*\query=\sw\to r,
\]
\[
\intro*\Prem=(\guarded{1},\link{1},\dots,\guarded{n},\link{n}),\qquad
\intro*\Mon{\Cons}{r}=\Prem\to(\query\to\Prem\to\goal)\to\goal .
\]
We call $\Mon{\Cons}{r}$ the \intro{monitor} of $\Cons$ and $r$. \AP Thus
$\Prem$ occurs twice; we call the members of the first occurrence
\intro{left} and those of the second \intro{right}, and write
$\guarded{i}^L,\guarded{i}^R$ and so on. Every \kl{occurrence} in a
\kl{left} member has, in $\Mon{\Cons}{r}$, the opposite \kl{polarity} of
the corresponding \kl{occurrence} in the \kl{right} member. \AP The
\intro{data variables} are the variables of $\Cons$ and $r$.

The \kl{switch} turns all constraints on at once, as needed to recover
the original formula; with the \kl{switch} off, any single constraint can be
turned on by its \kl{selector}, so that a countermodel for one modified
constraint need not satisfy the others. The \kl{right} copy of $\Prem$ is a
query that detects a modification of any \kl{left} member, and vice versa.

\begin{lemma}\label{lem:F-equiv}
Suppose that $\Cons\vdash_{\IL}\Phi\to r$, and that a substitution $\theta$
fixing the variables of $\Phi$ maps every $C_i$ to a theorem of $\IL$ and
$r$ to $\Phi$, where the \kl{selectors}, $\sw$ and $\goal$ do not occur in
$\Phi$. Then $\IL\vdash\Phi\to\Mon{\Cons}{r}$, and
$\IL\vdash\theta'(\Mon{\Cons}{r})\to\Phi$ for the extension $\theta'$ of
$\theta$ that maps $\sw$ and the \kl{selectors} to $\intro*\ttop=e\to e$ for a
new variable $e$ and $\goal$ to $\Phi$. Hence
$\IL\lplus\Mon{\Cons}{r}=\IL\lplus\Phi$.
\end{lemma}
\begin{proof}
Assume $\Phi$, the members of $\Prem$, and $\query\to\Prem\to\goal$.
Under the assumption $\sw$, the formulas $\link{i}$ and $\guarded{i}$ give
every $C_i$, hence $r$; thus $\query$ holds. Applying
$\query\to\Prem\to\goal$ to $\query$ and to the members of $\Prem$ again
gives $\goal$. Conversely, $\theta'$ maps every member of $\Prem$ to a
theorem of $\IL$ and $\query\to\Prem\to\goal$ to
$(\ttop\to\Phi)\to\theta'(\Prem)\to\Phi$, which is also a theorem.
\end{proof}

\begin{lemma}\label{lem:F-polar}
Every variable of $\Mon{\Cons}{r}$ other than $r$ has a \kl{positive} and a
\kl{negative} \kl{occurrence}; so does $r$ if it occurs in some $C_i$.
\end{lemma}
\begin{proof}
Every variable other than $\goal$ and $r$ occurs in $\Prem$, and
corresponding \kl{occurrences} in the \kl{left} and \kl{right} copies have
opposite \kl{polarities}. The variable $\goal$ occurs \kl{negatively} inside
$\query\to\Prem\to\goal$ and \kl{positively} at the end.
\end{proof}

\AP A formula is \intro{read-once} if no variable occurs in it twice.

\begin{lemma}\label{lem:readonce}
Let $\gamma$ be a \kl{read-once} formula without $\bot$, built by
$\wedge$, $\vee$ and $\to$, and let $u$ be a \kl{leaf} of $\gamma$. Then
there is a valuation of the variables other than the one at $u$ under which
the value of $\gamma$ is the value at $u$ or its complement; in particular
$\gamma$ takes both values.
\end{lemma}
\begin{proof}
By induction on $\gamma$, every such formula takes both values. Going from
$u$ to the root, give the sibling subformula the value $1$ at a conjunction
or at the antecedent of an implication whose consequent contains $u$, and
the value $0$ at a disjunction or at the consequent of an implication whose
antecedent contains $u$. Siblings have disjoint variables, so these
requirements are compatible.
\end{proof}

\begin{lemma}\label{lem:F-sep}
If every $C_i$ is \kl{read-once}, then no \kl{separation} of
$\Mon{\Cons}{r}$ belongs to $\CL$.
\end{lemma}
\begin{proof}
Let $\beta$ be obtained from $\Mon{\Cons}{r}$ by replacing one
\kl{occurrence} of a variable $q$ by a new variable $z$. Unless $q=\goal$,
we give $\goal$ the value $0$; then $\beta$ is false as soon as all
\kl{left} members are true and $\query$ or some \kl{right} member is false.
We give such a valuation in each case. By Lemma~\ref{lem:readonce}, for
every \kl{occurrence} in $C_i$ there are values of the \kl{data variables}
and of $z$ such that $C_i$ is false and the formula obtained from $C_i$ by
replacing that \kl{occurrence} by $z$ is true, and also values with the
opposite outcome.

\proofcase{(a) $q$ is a \kl{data variable} in $\guarded{i}^L$.}
Let $\sw=0$, $\sel{i}=1$ and $\sel{j}=0$ for $j\ne i$, and choose the
\kl{data variables} and $z$ so that $C_i$ is false and the modified
\kl{left} copy of $C_i$ is true. All \kl{left} members are true, and
$\guarded{i}^R$ is false.

\proofcase{(b) $q$ is a \kl{data variable} in $\guarded{i}^R$.}
As in (a), but with $C_i$ true and the modified \kl{right} copy false.

\proofcase{(c) $q=\sel{i}$ in $\guarded{i}^L$.}
Let $\sw=0$, $\sel{i}=1$, $z=0$, $\sel{j}=0$ for $j\ne i$, and make $C_i$
false. Then the \kl{left} members are true and $\guarded{i}^R$ is false.

\proofcase{(d) $q=\sel{i}$ in $\guarded{i}^R$.}
Let all \kl{selectors} and $\sw$ be $0$, $z=1$, and make $C_i$ false. Then
the modified $\guarded{i}^R$ is false.

\proofcase{(e) $q$ is $\sw$ or $\sel{i}$ in $\link{i}^L$.}
Let all \kl{data variables} be $1$, so that every $C_j$ is true. Let
$\sw=1$, $\sel{i}=0$ and $\sel{j}=1$ for $j\ne i$, and let $z=0$ if $q=\sw$
and $z=1$ if $q=\sel{i}$. The \kl{left} members are true, and
$\link{i}^R=\sw\to\sel{i}$ is false.

\proofcase{(f) $q=\sw$ in $\link{i}^R$, or in $\query$.}
Let $\sw$, all \kl{selectors} and $r$ be $0$, and $z=1$. The \kl{left}
members are true, and the modified $\link{i}^R$, resp.\ $\query$, is
false.

\proofcase{(g) $q=\sel{i}$ in $\link{i}^R$, or $q=r$ in $\query$.}
Let all variables other than $z$ and $\goal$ be $1$, and $z=0$. The
\kl{left} members are true, and the modified $\link{i}^R$, resp.\ $\query$,
is false.

\proofcase{(h) $q=\goal$.}
Let the \kl{negative} \kl{occurrence} of $\goal$ have value $1$, the
\kl{positive} one value $0$, and all other variables value $1$. All
antecedents of $\beta$ are true and its consequent is false.

Every \kl{occurrence} of a variable in $\Mon{\Cons}{r}$ is covered by one of
these cases.
\end{proof}

\begin{example}\label{ex:P-selectors}
For $\axP$ there are seven \kl{local constraints}, so $\Prem$ has fourteen
members, among them $\sel{001}\to b\to\nv{001}$ and $\sw\to\sel{001}$
(we index the \kl{selectors} by the \kl{occurrences} of $\axP$). Consider
again the \kl{separation} replacing $\nv{001}$ by $z$, now in the
\kl{left} copy of $\sel{001}\to b\to\nv{001}$. The valuation with $\sw=0$,
$\sel{001}=1$, all other \kl{selectors} $0$, $b=z=1$, $\nv{001}=0$ and
$\goal=0$ falsifies it: all \kl{left} members are true, since the
\kl{switch} and the other \kl{selectors} are off, whereas the \kl{right}
copy $\sel{001}\to b\to\nv{001}$ is false. No other constraint plays a
role.
\end{example}

\kl{Splittings} that modify several \kl{occurrences} at once may still
survive. For $\axP$, replace $\nv{001}$ by $z$ in \emph{both} copies of
$\sel{001}\to b\to\nv{001}$. The two copies remain identical, so the query
cannot detect the change; and under the \kl{switch} the modified constraint
$b\to z$ is not needed to obtain $\nv{\rt}$, as we saw above. Hence this
\kl{splitting} of $\Mon{\Cons_{\axP}}{\nv{\rt}}$ is in $\CL$, and so is the
\kl{abstraction} of both copies of $b\to\nv{001}$. The next construction
prevents modifications of this kind by making all \kl{occurrences}
distinct.

\subsection{The tree closure}\label{sec:minimality}

\AP We say that a formula $F$ is \intro{separation-rigid} if every variable
of $F$ has a \kl{positive} and a \kl{negative} \kl{occurrence} and no
\kl{separation} of $F$ belongs to $\CL$. We now show that a single further
construction turns a \kl{separation-rigid} formula of $\CL$ into a
\kl{minimal} one: we rename all \kl{occurrences} apart and tie the
\kl{ports} of each variable together by a spanning tree. A \kl{splitting}
of a \kl{center} then cuts an edge of the tree, and a \kl{splitting} of a
\kl{port} is either refuted by an edge or reduces to a \kl{separation} of
$F$.

Let $F$ be an implicational formula in which every variable has a
\kl{positive} and a \kl{negative} \kl{occurrence}, and let $\ported{F}$,
$\Guards_F$ be as in Section~\ref{sec:constraints}. \AP For a \kl{center}
$a$, let $\intro*\Pos{a}$ and $\intro*\Neg{a}$ be the sets of \kl{ports} of
\kl{positive} and \kl{negative} \kl{occurrences} of $a$, both nonempty.
Choose $p_a\in\Pos{a}$ and $n_a\in\Neg{a}$. The \intro{edge formulas} of
$a$ are
\[
n\to p_a\quad(n\in\Neg{a})
\qquad\text{and}\qquad
n_a\to p\quad(p\in\Pos{a}\setminus\{p_a\}).
\]
Viewed as edges, they form a spanning tree on $\Pos{a}\cup\Neg{a}$, and
every \kl{port} lies on at least one edge. \AP Let $\intro*\Edges_F$ be the
list of all \kl{edge formulas} and $\intro*\term$ a new variable. The
\intro{tree closure} of $F$ is
\[
\intro*\Tc{F}=\Guards_F\to(\ported{F}\to\Edges_F\to\term)\to\term .
\]
Since $\ported{F}\to\Edges_F\to\term$ is an antecedent of an antecedent,
every \kl{occurrence} in $\ported{F}$ has in $\Tc{F}$ the same
\kl{polarity} as the corresponding \kl{occurrence} in $F$.

\begin{example}\label{ex:P-ports}
In $\Mon{\Cons_{\axP}}{\nv{\rt}}$ the variable $\nv{001}$ occurs four times:
in the \kl{left} and \kl{right} copies of $\sel{00}\to\cons{00}$ and of
$\sel{001}\to\cons{001}$. Let $p_1,p_2,p_3,p_4$ be the \kl{ports} of these
\kl{occurrences}, in this order. The first and the last are \kl{positive}
and the other two \kl{negative}, so the \kl{guards} are
\[
\nv{001}\to p_1,\qquad p_2\to\nv{001},\qquad p_3\to\nv{001},\qquad
\nv{001}\to p_4,
\]
and, with $p_{\nv{001}}=p_1$ and $n_{\nv{001}}=p_2$, the \kl{edge formulas}
are $p_2\to p_1$, $p_3\to p_1$ and $p_2\to p_4$. The \kl{splitting}
described at the end of Section~\ref{sec:localization} becomes the
\kl{splitting} of the \kl{center} $\nv{001}$ that separates the \kl{guards}
of $p_3,p_4$ from those of $p_1,p_2$. The edge $p_3\to p_1$ joins the two
classes, and the valuation that gives the value $1$ to the \kl{center}
\kl{occurrence} and the \kl{port} on the side of $p_3$, and $0$ on the side
of $p_1$, makes all \kl{guards} true and $p_3\to p_1$ false. The
\kl{abstraction} of $b\to\nv{001}$ disappears as well, since its two
copies become $p\to p_3$ and $p'\to p_4$ for distinct \kl{ports} $p,p'$ of
$b$.
\end{example}

\begin{lemma}\label{lem:M-equiv}
$\IL\vdash F\to\Tc{F}$, and $\IL\vdash\collsub(\Tc{F})\to F$ for the
extension $\collsub$ of $\collsub_F$ that maps $\term$ to $F$. Hence
$\IL\lplus\Tc{F}=\IL\lplus F$.
\end{lemma}
\begin{proof}
By Lemma~\ref{lem:port}, $\Guards_F\vdash_{\IL}F\to\ported{F}$, and every
\kl{edge formula} $n\to p$ of $a$ follows from the \kl{guards} $n\to a$ and
$a\to p$. Hence from $F$, the \kl{guards} and $\ported{F}\to\Edges_F\to\term$
we obtain $\term$. Conversely, $\collsub$ maps every \kl{guard} and every
\kl{edge formula} to a formula $a\to a$, and maps
$\ported{F}\to\Edges_F\to\term$ to $F\to\collsub(\Edges_F)\to F$.
\end{proof}

\begin{lemma}\label{lem:M-split}
If $F$ is \kl{separation-rigid}, then no \kl{splitting} of $\Tc{F}$ belongs
to $\CL$.
\end{lemma}
\begin{proof}
Let $\beta$ be a \kl{splitting} of $\Tc{F}$ that replaces the
\kl{occurrences} in a set $Q$ of a variable $q$ by $z$. Unless $q=\term$,
we give $\term$ the value $0$; then $\beta$ is false as soon as all
\kl{guards} are true and $\ported{F}$ or some \kl{edge formula} is false.

\proofcase{(a) $q$ is a \kl{center} $a$.}
The variable $a$ occurs exactly once in each \kl{guard} of an
\kl{occurrence} of $a$ and nowhere else, so $Q$ divides the \kl{ports} of
$a$ into two nonempty classes. Since the \kl{edge formulas} of $a$ form a
spanning tree, some \kl{edge formula} $n\to p$ of $a$ joins the classes. Give
the \kl{occurrence} of $a$ or $z$ in the \kl{guard} of $n$ the value $1$
and the one in the \kl{guard} of $p$ the value $0$, give each \kl{port} of
$a$ the value of the \kl{center} occurrence in its \kl{guard}, and give all
other variables the value $0$. All \kl{guards} are true and $n\to p$ is
false.

\proofcase{(b) $q$ is a \kl{port} $\port{o}$ and some \kl{occurrence} of
$\port{o}$ in an \kl{edge formula} lies on the other side from the one in
its \kl{guard}.}
We may assume that the \kl{occurrence} in the \kl{guard} is not in $Q$
(otherwise interchange $q$ and $z$, which yields a \kl{renaming} of
$\beta$). If $\port{o}$ is \kl{positive}, let $z=0$ and all other variables
$1$; then all \kl{guards} are true, and an \kl{edge formula} with $z$ as
its consequent is false. If $\port{o}$ is \kl{negative}, let $z=1$ and all
other variables $0$; then an \kl{edge formula} with $z$ as its antecedent is
false.

\proofcase{(c) $q$ is a \kl{port} $\port{o}$ and all its \kl{occurrences} in
\kl{edge formulas} lie on the side of its \kl{guard}.}
Since $\port{o}$ occurs exactly once in $\ported{F}$ and $Q$ is a proper
nonempty subset, we may assume that $Q$ is the \kl{occurrence} in
$\ported{F}$. Since $F$ is \kl{separation-rigid}, the \kl{separation} of
$F$ at $o$ has a falsifying valuation $v$. Give every \kl{center} its value
under $v$, every \kl{port} except $z$ the value of its \kl{center}, and $z$
the value of the new variable of the \kl{separation}. All \kl{guards} and
\kl{edge formulas} are true, and the modified $\ported{F}$ has the value of
the \kl{separation} of $F$ under $v$, that is, $0$.

\proofcase{(d) $q=\term$.}
Give the inner \kl{occurrence} of $\term$ the value $1$, the outer one $0$,
and all other variables $1$.
\end{proof}

\begin{lemma}\label{lem:abs-necessary}
Let $\alpha\in\CL$ be a formula no \kl{splitting} of which belongs to
$\CL$, and suppose that an \kl{abstraction} $\gen{\alpha}{Q}{z}$, where $Q$
is a set of \kl{occurrences} of a compound formula $\gamma$, belongs to
$\CL$. Then there are distinct variables $a,b$ such that every
\kl{occurrence} of $a$ and of $b$ in $\alpha$ lies in an \kl{occurrence} of
$a\to b$.
\end{lemma}
\begin{proof}
Let $\ell$ be a \kl{leaf} of $\gamma$, with variable $a$, and let $\gamma'$
be $\gamma$ with $a$ replaced by a new variable $a'$ at $\ell$. Substituting
$\gamma'$ for $z$ in $\gamma[Q:=z]$ gives a formula of $\CL$, namely the
formula obtained from $\alpha$ by replacing $a$ by $a'$ at the set
$S_\ell=\{u\ell\mid u\in Q\}$. Since this is not a \kl{splitting},
$S_\ell=\Occ_\alpha(a)$. If $a$ also occurred at another \kl{leaf} $k$ of
$\gamma$, then $S_k=\Occ_\alpha(a)$ as well, although $S_k\cap
S_\ell=\emptyset$. Hence $\gamma$ is \kl{read-once} and every
\kl{occurrence} of each of its variables lies at the corresponding position
of an element of $Q$. Taking an \kl{occurrence} in $\gamma$ of a
subformula $a\to b$ with $a,b$ variables, we obtain the claim.
\end{proof}

\begin{lemma}\label{lem:M-noabs}
In $\Tc{F}$ there are no distinct variables $a,b$ such that every
\kl{occurrence} of $a$ and of $b$ lies in an \kl{occurrence} of $a\to b$.
\end{lemma}
\begin{proof}
A subformula of $\Tc{F}$ of the form $a\to b$ with $a,b$ variables is a
\kl{guard}, or occurs in $\ported{F}$, or is an \kl{edge formula}; the
antecedents of the outer implications of $\Tc{F}$ and of
$\ported{F}\to\Edges_F\to\term$ are compound, and $\term$ is not the
consequent of an implication with a variable antecedent. A \kl{guard}
contains a \kl{port}, which also occurs in $\ported{F}$. The other two
kinds consist of \kl{ports}, each of which also occurs in its \kl{guard},
which contains a \kl{center}.
\end{proof}

\begin{theorem}\label{thm:closure}
Let $F\in\CL$ be \kl{separation-rigid}. Then $\Tc{F}$ is \kl{minimal in}
$\CL$, $\IL\lplus\Tc{F}=\IL\lplus F$, $\Ord(\Tc{F})\le\Ord(F)+2$, and
$\sizeof{\Tc{F}}=5\sizeof{F}-2|\Var(F)|+2$.
\end{theorem}
\begin{proof}
We have $\Tc{F}\in\CL$ by Lemma~\ref{lem:M-equiv}. No \kl{splitting} of
$\Tc{F}$ is in $\CL$ by Lemma~\ref{lem:M-split}, and hence no
\kl{abstraction} by Lemmas~\ref{lem:abs-necessary} and~\ref{lem:M-noabs}.
By Proposition~\ref{prop:criterion}, $\Tc{F}$ is \kl{minimal in} $\CL$.
\kl{Guards} and \kl{edge formulas} have \kl{order} $2$, which gives the
bound on the \kl{order}. There are $\sizeof{F}$ \kl{guards} and
$\sizeof{F}-|\Var(F)|$ \kl{edge formulas}, which gives the \kl{size}.
\end{proof}

\subsection{Minimal axioms}\label{sec:first-layer}

Let $\Phi$ be an implicational formula that is not a variable, and put
\[
\intro*\Fst{\Phi}=\Mon{\Cons_\Phi}{\nv{\rt}},\qquad
\intro*\Mz{\Phi}=\Tc{\Fst{\Phi}} .
\]

\begin{theorem}\label{thm:first-layer}
Let $\Phi\in\CL$ be an implicational formula that is not a variable. Then
$\Mz{\Phi}$ is \kl{minimal in} $\CL$, and for the substitution
$\intro*\recsub{\Phi}=\theta'\circ\collsub$, which fixes the variables of
$\Phi$, where $\theta'$ is as in Lemma~\ref{lem:F-equiv} for the
substitution of Lemma~\ref{lem:constraints} and $\collsub$ is as in
Lemma~\ref{lem:M-equiv},
\[
\IL\vdash\Phi\to\Mz{\Phi},\qquad
\IL\vdash\recsub{\Phi}(\Mz{\Phi})\to\Phi,\qquad
\IL\lplus\Mz{\Phi}=\IL\lplus\Phi .
\]
Moreover
$\sizeof{\Mz{\Phi}}=102\sizeof{\Phi}-2|\Var(\Phi)|-38$ and
$\Ord(\Mz{\Phi})\le7$.
\end{theorem}
\begin{proof}
By Lemmas~\ref{lem:constraints} and~\ref{lem:F-equiv},
$\IL\vdash\Phi\to\Fst{\Phi}$ and $\IL\vdash\theta'(\Fst{\Phi})\to\Phi$; in
particular $\Fst{\Phi}\in\CL$. The \kl{local constraints} are
\kl{read-once} and $\nv{\rt}$ occurs in $\cons{\rt}$, so $\Fst{\Phi}$ is
\kl{separation-rigid} by Lemmas~\ref{lem:F-polar} and~\ref{lem:F-sep}.
Theorem~\ref{thm:closure} and Lemma~\ref{lem:M-equiv} give the
\kl{minimality} and the equivalences.

For the size, let $L=\sizeof{\Phi}$ and $V=|\Var(\Phi)|$. The $2L-1$
\kl{local constraints} have total size $2L+3(L-1)=5L-3$, so the members of
$\Prem$ have total size $5L-3+3(2L-1)=11L-6$, and
$\sizeof{\Fst{\Phi}}=2(11L-6)+4=22L-8$. The variables of $\Fst{\Phi}$ are
those of $\Phi$, the $2L-1$ \kl{node variables}, the $2L-1$
\kl{selectors}, $\sw$ and $\goal$, that is, $4L+V$ in number.
Theorem~\ref{thm:closure} gives the formula. Finally, the \kl{local
constraints} have \kl{order} at most $3$, so $\Ord(\Fst{\Phi})\le5$ and
$\Ord(\Mz{\Phi})\le7$.
\end{proof}
